% Options for packages loaded elsewhere
% Options for packages loaded elsewhere
\PassOptionsToPackage{unicode}{hyperref}
\PassOptionsToPackage{hyphens}{url}
%
\documentclass[
  ignorenonframetext,
  aspectratio=169,
  spanish,
]{beamer}
\newif\ifbibliography
\usepackage{pgfpages}
\setbeamertemplate{caption}[numbered]
\setbeamertemplate{caption label separator}{: }
\setbeamercolor{caption name}{fg=normal text.fg}
\beamertemplatenavigationsymbolshorizontal
% remove section numbering
\setbeamertemplate{part page}{
  \centering
  \begin{beamercolorbox}[sep=16pt,center]{part title}
    \usebeamerfont{part title}\insertpart\par
  \end{beamercolorbox}
}
\setbeamertemplate{section page}{
  \centering
  \begin{beamercolorbox}[sep=12pt,center]{section title}
    \usebeamerfont{section title}\insertsection\par
  \end{beamercolorbox}
}
\setbeamertemplate{subsection page}{
  \centering
  \begin{beamercolorbox}[sep=8pt,center]{subsection title}
    \usebeamerfont{subsection title}\insertsubsection\par
  \end{beamercolorbox}
}
% Prevent slide breaks in the middle of a paragraph
\widowpenalties 1 10000
\raggedbottom
\AtBeginPart{
  \frame{\partpage}
}
\AtBeginSection{
  \ifbibliography
  \else
    \frame{\sectionpage}
  \fi
}
\AtBeginSubsection{
  \frame{\subsectionpage}
}
\usepackage{iftex}
\ifPDFTeX
  \usepackage[T1]{fontenc}
  \usepackage[utf8]{inputenc}
  \usepackage{textcomp} % provide euro and other symbols
\else % if luatex or xetex
  \usepackage{unicode-math} % this also loads fontspec
  \defaultfontfeatures{Scale=MatchLowercase}
  \defaultfontfeatures[\rmfamily]{Ligatures=TeX,Scale=1}
\fi
\usepackage{lmodern}

\usetheme[]{default}
\usefonttheme[]{default}
\ifPDFTeX\else
  % xetex/luatex font selection
\fi
% Use upquote if available, for straight quotes in verbatim environments
\IfFileExists{upquote.sty}{\usepackage{upquote}}{}
\IfFileExists{microtype.sty}{% use microtype if available
  \usepackage[]{microtype}
  \UseMicrotypeSet[protrusion]{basicmath} % disable protrusion for tt fonts
}{}
\makeatletter
\@ifundefined{KOMAClassName}{% if non-KOMA class
  \IfFileExists{parskip.sty}{%
    \usepackage{parskip}
  }{% else
    \setlength{\parindent}{0pt}
    \setlength{\parskip}{6pt plus 2pt minus 1pt}}
}{% if KOMA class
  \KOMAoptions{parskip=half}}
\makeatother

\usepackage{color}
\usepackage{fancyvrb}
\newcommand{\VerbBar}{|}
\newcommand{\VERB}{\Verb[commandchars=\\\{\}]}
\DefineVerbatimEnvironment{Highlighting}{Verbatim}{commandchars=\\\{\}}
% Add ',fontsize=\small' for more characters per line
\usepackage{framed}
\definecolor{shadecolor}{RGB}{241,243,245}
\newenvironment{Shaded}{\begin{snugshade}}{\end{snugshade}}
\newcommand{\AlertTok}[1]{\textcolor[rgb]{0.68,0.00,0.00}{#1}}
\newcommand{\AnnotationTok}[1]{\textcolor[rgb]{0.37,0.37,0.37}{#1}}
\newcommand{\AttributeTok}[1]{\textcolor[rgb]{0.40,0.45,0.13}{#1}}
\newcommand{\BaseNTok}[1]{\textcolor[rgb]{0.68,0.00,0.00}{#1}}
\newcommand{\BuiltInTok}[1]{\textcolor[rgb]{0.00,0.23,0.31}{#1}}
\newcommand{\CharTok}[1]{\textcolor[rgb]{0.13,0.47,0.30}{#1}}
\newcommand{\CommentTok}[1]{\textcolor[rgb]{0.37,0.37,0.37}{#1}}
\newcommand{\CommentVarTok}[1]{\textcolor[rgb]{0.37,0.37,0.37}{\textit{#1}}}
\newcommand{\ConstantTok}[1]{\textcolor[rgb]{0.56,0.35,0.01}{#1}}
\newcommand{\ControlFlowTok}[1]{\textcolor[rgb]{0.00,0.23,0.31}{\textbf{#1}}}
\newcommand{\DataTypeTok}[1]{\textcolor[rgb]{0.68,0.00,0.00}{#1}}
\newcommand{\DecValTok}[1]{\textcolor[rgb]{0.68,0.00,0.00}{#1}}
\newcommand{\DocumentationTok}[1]{\textcolor[rgb]{0.37,0.37,0.37}{\textit{#1}}}
\newcommand{\ErrorTok}[1]{\textcolor[rgb]{0.68,0.00,0.00}{#1}}
\newcommand{\ExtensionTok}[1]{\textcolor[rgb]{0.00,0.23,0.31}{#1}}
\newcommand{\FloatTok}[1]{\textcolor[rgb]{0.68,0.00,0.00}{#1}}
\newcommand{\FunctionTok}[1]{\textcolor[rgb]{0.28,0.35,0.67}{#1}}
\newcommand{\ImportTok}[1]{\textcolor[rgb]{0.00,0.46,0.62}{#1}}
\newcommand{\InformationTok}[1]{\textcolor[rgb]{0.37,0.37,0.37}{#1}}
\newcommand{\KeywordTok}[1]{\textcolor[rgb]{0.00,0.23,0.31}{\textbf{#1}}}
\newcommand{\NormalTok}[1]{\textcolor[rgb]{0.00,0.23,0.31}{#1}}
\newcommand{\OperatorTok}[1]{\textcolor[rgb]{0.37,0.37,0.37}{#1}}
\newcommand{\OtherTok}[1]{\textcolor[rgb]{0.00,0.23,0.31}{#1}}
\newcommand{\PreprocessorTok}[1]{\textcolor[rgb]{0.68,0.00,0.00}{#1}}
\newcommand{\RegionMarkerTok}[1]{\textcolor[rgb]{0.00,0.23,0.31}{#1}}
\newcommand{\SpecialCharTok}[1]{\textcolor[rgb]{0.37,0.37,0.37}{#1}}
\newcommand{\SpecialStringTok}[1]{\textcolor[rgb]{0.13,0.47,0.30}{#1}}
\newcommand{\StringTok}[1]{\textcolor[rgb]{0.13,0.47,0.30}{#1}}
\newcommand{\VariableTok}[1]{\textcolor[rgb]{0.07,0.07,0.07}{#1}}
\newcommand{\VerbatimStringTok}[1]{\textcolor[rgb]{0.13,0.47,0.30}{#1}}
\newcommand{\WarningTok}[1]{\textcolor[rgb]{0.37,0.37,0.37}{\textit{#1}}}

\usepackage{longtable,booktabs,array}
\usepackage{calc} % for calculating minipage widths
\usepackage{caption}
% Make caption package work with longtable
\makeatletter
\def\fnum@table{\tablename~\thetable}
\makeatother
\usepackage{graphicx}
\makeatletter
\newsavebox\pandoc@box
\newcommand*\pandocbounded[1]{% scales image to fit in text height/width
  \sbox\pandoc@box{#1}%
  \Gscale@div\@tempa{\textheight}{\dimexpr\ht\pandoc@box+\dp\pandoc@box\relax}%
  \Gscale@div\@tempb{\linewidth}{\wd\pandoc@box}%
  \ifdim\@tempb\p@<\@tempa\p@\let\@tempa\@tempb\fi% select the smaller of both
  \ifdim\@tempa\p@<\p@\scalebox{\@tempa}{\usebox\pandoc@box}%
  \else\usebox{\pandoc@box}%
  \fi%
}
% Set default figure placement to htbp
\def\fps@figure{htbp}
\makeatother



\ifLuaTeX
\usepackage[bidi=basic]{babel}
\else
\usepackage[bidi=default]{babel}
\fi
% get rid of language-specific shorthands (see #6817):
\let\LanguageShortHands\languageshorthands
\def\languageshorthands#1{}


\setlength{\emergencystretch}{3em} % prevent overfull lines

\providecommand{\tightlist}{%
  \setlength{\itemsep}{0pt}\setlength{\parskip}{0pt}}



 


\setbeamertemplate{headline}{} % Elimina la cabecera del tema Madrid
\setbeamertemplate{section in head/foot}{} % Limpia la numeración de secciones
\renewcommand{\contentsname}{Índice}
\renewcommand{\refname}{Referencias}
\renewcommand{\figurename}{Figura}
\renewcommand{\tablename}{Tabla}
\logo{%
  \begin{tabular}{c}
    \includegraphics[height=1.2cm]{EscudoU.png}
  \end{tabular}%
}
% --- CONFIGURACIÓN DE COLORES INSTITUCIONALES ---
\definecolor{UnicorGreen}{HTML}{024426}
\definecolor{UnicorText}{HTML}{F9F9F9}
\definecolor{UnicorWhite}{HTML}{FFFFFF}
\definecolor{UnicorGold}{HTML}{D1B061}

% Aplicación al fondo y texto general
\setbeamercolor{background canvas}{bg=UnicorGreen}
\setbeamercolor{normal text}{fg=UnicorText}

% Aplicación a títulos y cabeceras
%\setbeamercolor{title}{fg=UnicorWhite}
\setbeamercolor{frametitle}{fg=UnicorWhite}
\setbeamercolor{framesubtitle}{fg=UnicorWhite}

% Aplicación a viñetas y estructura
\setbeamercolor{structure}{fg=UnicorWhite}
\setbeamercolor{itemize item}{fg=UnicorWhite}
\setbeamercolor{itemize subitem}{fg=UnicorWhite}
\setbeamercolor{enumerate item}{fg=UnicorWhite}

% Forzar a LaTeX a repintar el texto general al iniciar el documento
\AtBeginDocument{\usebeamercolor[fg]{normal text}}

% Aplicación a hipervínculos
\hypersetup{colorlinks=true, urlcolor=UnicorGold, linkcolor=UnicorWhite}

% --- ACTIVAR ESTILO DE BLOQUES Y SOMBRAS (ESTILO MADRID) ---
\useinnertheme[shadow=true]{rounded}

% --- COLORES DEL RECUADRO DEL TÍTULO PRINCIPAL ---
% Fondo dorado con texto verde
\setbeamercolor{title}{bg=UnicorGold, fg=UnicorGreen}
\setbeamercolor{frametitle}{fg=UnicorWhite} % Mantiene el título de las diapositivas normal
\setbeamercolor{framesubtitle}{fg=UnicorWhite}

% --- COLORES PARA TEOREMAS, DEFINICIONES Y BLOQUES ---
% Título del bloque (Fondo dorado, texto verde)
\setbeamercolor{block title}{bg=UnicorGold, fg=UnicorGreen}
% Cuerpo del bloque (Fondo blanco, texto verde)
\setbeamercolor{block body}{bg=UnicorWhite, fg=UnicorGreen}

\AtBeginDocument{
  \author{%
    \begin{tabular}{c}
      Mario A. Morales Rivera \\
      \footnotesize \texttt{mamorales@correo.unicordoba.edu.co} \\
      \footnotesize Universidad de Córdoba
    \end{tabular}
  }
  \institute{}
}
% --- LIMPIAR DIAPOSITIVAS DE SECCIÓN (ELIMINAR "SECTION X") ---
\setbeamertemplate{section page}{
  \begin{center}
    \begin{beamercolorbox}[sep=12pt,center,rounded=true,shadow=true]{title}
      \usebeamerfont{title}\insertsection\par
    \end{beamercolorbox}
  \end{center}
}
\makeatletter
\@ifpackageloaded{tcolorbox}{}{\usepackage[skins,breakable]{tcolorbox}}
\@ifpackageloaded{fontawesome5}{}{\usepackage{fontawesome5}}
\definecolor{quarto-callout-color}{HTML}{909090}
\definecolor{quarto-callout-note-color}{HTML}{0758E5}
\definecolor{quarto-callout-important-color}{HTML}{CC1914}
\definecolor{quarto-callout-warning-color}{HTML}{EB9113}
\definecolor{quarto-callout-tip-color}{HTML}{00A047}
\definecolor{quarto-callout-caution-color}{HTML}{FC5300}
\definecolor{quarto-callout-color-frame}{HTML}{acacac}
\definecolor{quarto-callout-note-color-frame}{HTML}{4582ec}
\definecolor{quarto-callout-important-color-frame}{HTML}{d9534f}
\definecolor{quarto-callout-warning-color-frame}{HTML}{f0ad4e}
\definecolor{quarto-callout-tip-color-frame}{HTML}{02b875}
\definecolor{quarto-callout-caution-color-frame}{HTML}{fd7e14}
\makeatother
\makeatletter
\@ifpackageloaded{caption}{}{\usepackage{caption}}
\AtBeginDocument{%
\ifdefined\contentsname
  \renewcommand*\contentsname{Tabla de contenidos}
\else
  \newcommand\contentsname{Tabla de contenidos}
\fi
\ifdefined\listfigurename
  \renewcommand*\listfigurename{Listado de Figuras}
\else
  \newcommand\listfigurename{Listado de Figuras}
\fi
\ifdefined\listtablename
  \renewcommand*\listtablename{Listado de Tablas}
\else
  \newcommand\listtablename{Listado de Tablas}
\fi
\ifdefined\figurename
  \renewcommand*\figurename{Figura}
\else
  \newcommand\figurename{Figura}
\fi
\ifdefined\tablename
  \renewcommand*\tablename{Tabla}
\else
  \newcommand\tablename{Tabla}
\fi
}
\@ifpackageloaded{float}{}{\usepackage{float}}
\floatstyle{ruled}
\@ifundefined{c@chapter}{\newfloat{codelisting}{h}{lop}}{\newfloat{codelisting}{h}{lop}[chapter]}
\floatname{codelisting}{Listado}
\newcommand*\listoflistings{\listof{codelisting}{Listado de Listados}}
\makeatother
\makeatletter
\makeatother
\makeatletter
\@ifpackageloaded{caption}{}{\usepackage{caption}}
\@ifpackageloaded{subcaption}{}{\usepackage{subcaption}}
\makeatother
\makeatletter
\@ifpackageloaded{tikz}{}{\usepackage{tikz}}
\makeatother
        \newcommand*\circled[1]{\tikz[baseline=(char.base)]{
          \node[shape=circle,draw,inner sep=1pt] (char) {{\scriptsize#1}};}}  
                  

\usepackage{bookmark}
\IfFileExists{xurl.sty}{\usepackage{xurl}}{} % add URL line breaks if available
\urlstyle{same}
\hypersetup{
  pdftitle={Diccionarios},
  pdfauthor={Mario Morales},
  pdflang={es},
  hidelinks,
  pdfcreator={LaTeX via pandoc}}


\title{Diccionarios}
\subtitle{Estadística con Python}
\author{Mario Morales}
\date{}

\begin{document}
\frame{\titlepage}

\renewcommand*\contentsname{Tabla de contenidos}
\begin{frame}[allowframebreaks]
  \frametitle{Tabla de contenidos}
  \setcounter{tocdepth}{3}
  \tableofcontents
\end{frame}
\setcounter{tocdepth}{3}
\tableofcontents
}

\section{Introducción}\label{introducciuxf3n}

\begin{frame}{Diccionarios}
\phantomsection\label{diccionarios}
\index{Diccionarios}

\begin{tcolorbox}[enhanced jigsaw, leftrule=.75mm, titlerule=0mm, bottomrule=.15mm, opacitybacktitle=0.6, colbacktitle=quarto-callout-note-color!10!white, arc=.35mm, colframe=quarto-callout-note-color-frame, bottomtitle=1mm, title=\textcolor{quarto-callout-note-color}{\faInfo}\hspace{0.5em}{Nota}, coltitle=black, colback=white, toptitle=1mm, rightrule=.15mm, breakable, toprule=.15mm, left=2mm, opacityback=0]

Un diccionario es un objeto muy flexible de Python para almacenar
diferentes tipos de datos y objetos, tales como números, caracteres,
listas y arreglos.

\end{tcolorbox}

\begin{itemize}
\item
  Como estudiamos en clase anterior, una lista es una colección de
  objetos indexada desde cero (0) hasta el número de objetos menos uno
  \((n-1)\).
\item
  En lugar de buscar un elemento a través de un índice de números
  enteros puede ser más fácil, práctico (o útil) usar una cadena de
  texto con un nombre, que sea fácil de recordar.
\end{itemize}
\end{frame}

\begin{frame}[fragile]{Diccionarios (cont)}
\phantomsection\label{diccionarios-cont}
\begin{itemize}
\item
  Los diccionarios nos permiten identificar cada elemento por una clave
  \texttt{(key)} que por lo general es un \emph{string} o cadena de
  caracteres.
\item
  Rigurosamente hablando, un diccionario en Python es una lista donde
  los índices pueden ser cadenas de texto.
\item
  Suponga que necesitamos guardar las edades de tres personas: Brenda,
  Luis y Diego. Para eso podemos usar una lista
\end{itemize}

\begin{Shaded}
\begin{Highlighting}[]
\NormalTok{edad}\OperatorTok{=}\NormalTok{[}\DecValTok{22}\NormalTok{,}\DecValTok{15}\NormalTok{,}\DecValTok{19}\NormalTok{]  }
\end{Highlighting}
\end{Shaded}
\end{frame}

\begin{frame}[fragile]{Diccionarios (cont)}
\phantomsection\label{diccionarios-cont-1}
\begin{itemize}
\item
  Pero de esta forma debemos recordar la secuencia de nombres, por
  ejemplo que el índice cero corresponde a Brenda, el índice 1 a Luis y
  el 2 a Diego.
\item
  Esto es, la edad de Luis se obtiene como \texttt{edad{[}1{]}}. Un
  diccionario con los nombres de las personas como índice es más
  conveniente porque esto permite escribir
  \texttt{edad{[}\textquotesingle{}Luis\textquotesingle{}{]}} para
  recuperar la edad de Luis, sin tener que recordar el índice numérico.
\end{itemize}
\end{frame}

\section{Creación de diccionarios}\label{creaciuxf3n-de-diccionarios}

\begin{frame}[fragile]{Creación de diccionarios}
\phantomsection\label{sec-crearciondiccionarios}
\index{Diccionarios} \index{Diccionarios!Creación}

En esta sección se estudian las formas de crear diccionarios, primero
mediante \texttt{\{clave:valor,...\}} luego mediante la función
\texttt{dict}
\index{Diccionarios!Creación!mediante \texttt{\{clave:valor\}} }

A continuación se crea un diccionario que contiene los nombres de las
personas como claves y sus edades como valores

\begin{Shaded}
\begin{Highlighting}[]
\NormalTok{edad}\OperatorTok{=}\NormalTok{\{}\StringTok{\textquotesingle{}Luis\textquotesingle{}}\NormalTok{:}\DecValTok{15}\NormalTok{,}\StringTok{\textquotesingle{}Brenda\textquotesingle{}}\NormalTok{:}\DecValTok{22}\NormalTok{,}\StringTok{\textquotesingle{}Diego\textquotesingle{}}\NormalTok{:}\DecValTok{19}\NormalTok{\}}
\end{Highlighting}
\end{Shaded}

\index{Diccionarios!Creación!mediante \texttt{dict()}}

\begin{block}{Creación de diccionarios: función \texttt{dict()}}
\phantomsection\label{creaciuxf3n-de-diccionarios-funciuxf3n-dict}
De manera equivalente, podemos usar la función \texttt{dict()} para
crear el mismo diccionario. En este caso entregamos los datos de la
forma \texttt{clave=valor} como se muestra a continuación

\begin{Shaded}
\begin{Highlighting}[]
\NormalTok{edad}\OperatorTok{=}\BuiltInTok{dict}\NormalTok{(Luis}\OperatorTok{=}\DecValTok{15}\NormalTok{,Brenda}\OperatorTok{=}\DecValTok{22}\NormalTok{,Diego}\OperatorTok{=}\DecValTok{19}\NormalTok{) }
\end{Highlighting}
\end{Shaded}
\end{block}

\begin{block}{Creación de diccionarios: \texttt{dict(zip())}}
\phantomsection\label{creaciuxf3n-de-diccionarios-dictzip}
\index{Diccionarios!Creación!mediante \texttt{dict(zip())}}

\begin{itemize}
\item
  Suponga que tenemos los nombres de las personas en una lista o tupla y
  sus edades en otra lista o tupla, podemos crear el diccionario
  mediante \texttt{dict(zip(a,b))}.
\item
  En este caso se le entrega a la función \texttt{zip()} dos objetos
  iterables, \texttt{a} y \texttt{b}, de la misma longitud, que podrían
  ser una cadena de caracteres, una lista o una tupla,
\item
  \texttt{dict()} devolverá un diccionario cuya clave \(i-\)ésima es el
  elemento \(i\) del primer iterable (\texttt{a}) y el valor asignado a
  esa clave es \(i-\)ésimo del segundo iterable (\texttt{b})
\end{itemize}
\end{block}

\begin{block}{Creación de diccionarios: \texttt{dict(zip())}, ejemplo}
\phantomsection\label{creaciuxf3n-de-diccionarios-dictzip-ejemplo}
\phantomsection\label{annotated-cell-4}%
\begin{Shaded}
\begin{Highlighting}[]
\NormalTok{nombres}\OperatorTok{=}\NormalTok{(}\StringTok{"Luis"}\NormalTok{,}\StringTok{"Brenda"}\NormalTok{,}\StringTok{"Diego"}\NormalTok{) }\hspace*{\fill}\NormalTok{\circled{1}}
\NormalTok{edades}\OperatorTok{=}\NormalTok{[}\DecValTok{15}\NormalTok{,}\DecValTok{22}\NormalTok{,}\DecValTok{19}\NormalTok{] }\hspace*{\fill}\NormalTok{\circled{2}}
\NormalTok{edad}\OperatorTok{=}\BuiltInTok{dict}\NormalTok{(}\BuiltInTok{zip}\NormalTok{(nombres,edades)) }\hspace*{\fill}\NormalTok{\circled{3}}
\end{Highlighting}
\end{Shaded}

\begin{description}
\tightlist
\item[\circled{1}]
Tupla que contiene los nombres.
\item[\circled{2}]
Lista que contiene las correspondientes edades.
\item[\circled{3}]
Diccionario con los nombres como como claves y las edades como valores.
\end{description}
\end{block}

\begin{block}{Diccionarios: Claves y valores}
\phantomsection\label{diccionarios-claves-y-valores}
\begin{itemize}
\item
  Los índices de texto en los diccionarios se conocen como claves
  \emph{(Keys)}.
\item
  En el diccionario \texttt{edad}, creado mediante los ejemplos
  anteriores, las claves son
  \texttt{\textquotesingle{}Luis\textquotesingle{}},
  \texttt{\textquotesingle{}Brenda\textquotesingle{}} y
  \texttt{\textquotesingle{}Diego\textquotesingle{}} y los valores son
  15, 22 y 19.
\end{itemize}
\end{block}
\end{frame}

\section{Edición de diccionarios}\label{ediciuxf3n-de-diccionarios}

\begin{frame}[fragile]{Adición de elementos a un diccionario}
\phantomsection\label{adiciuxf3n-de-elementos-a-un-diccionario}
En cualquier momento, después de su creación, se puede adicionar una
pareja \texttt{clave:valor} a un diccionario, por ejemplo, a partir del
diccionario \texttt{edad} agregamos a \texttt{Jaime} , cuya edad es
\texttt{21} años

\begin{Shaded}
\begin{Highlighting}[]
\NormalTok{edad[}\StringTok{"Jaime"}\NormalTok{]}\OperatorTok{=}\DecValTok{21} 
\BuiltInTok{print}\NormalTok{(edad)}
\end{Highlighting}
\end{Shaded}

\begin{verbatim}
{'Luis': 15, 'Brenda': 22, 'Diego': 19, 'Jaime': 21}
\end{verbatim}
\end{frame}

\begin{frame}[fragile]{Diccionarios: tipos de valores}
\phantomsection\label{diccionarios-tipos-de-valores}
En el ejemplo anterior todos los valores del diccionario son números
enteros, pero es posible crear diccionarios con diferentes tipos de
valores, como en el siguiente ejemplo

\begin{Shaded}
\begin{Highlighting}[]
\NormalTok{d1}\OperatorTok{=}\NormalTok{ \{}
     \StringTok{"clave1"}\NormalTok{:}\DecValTok{234}\NormalTok{,}
     \StringTok{"clave2"}\NormalTok{:}\VariableTok{True}\NormalTok{,}
     \StringTok{"clave3"}\NormalTok{:}\StringTok{"Valor 1"}\NormalTok{,}
     \StringTok{"clave4"}\NormalTok{:[}\DecValTok{1}\NormalTok{,}\DecValTok{2}\NormalTok{,}\DecValTok{3}\NormalTok{,}\DecValTok{4}\NormalTok{]\}}
\end{Highlighting}
\end{Shaded}
\end{frame}

\begin{frame}[fragile]{Diccionarios: tipos de valores (II)}
\phantomsection\label{diccionarios-tipos-de-valores-ii}
como con los objetos estudiados en las secciones anteriores, podemos
usar la la función \texttt{type()} para confirmar si un objeto de Python
es un diccionario y obtener el tipo de datos que tienen sus elementos

\phantomsection\label{annotated-cell-7}%
\begin{Shaded}
\begin{Highlighting}[]
\BuiltInTok{type}\NormalTok{(d1) }\hspace*{\fill}\NormalTok{\circled{1}}
\BuiltInTok{type}\NormalTok{(d1) }\KeywordTok{is} \BuiltInTok{dict} \hspace*{\fill}\NormalTok{\circled{2}}
\BuiltInTok{type}\NormalTok{(d1[}\StringTok{\textquotesingle{}clave2\textquotesingle{}}\NormalTok{]) }\KeywordTok{is} \BuiltInTok{int} \hspace*{\fill}\NormalTok{\circled{3}}
\end{Highlighting}
\end{Shaded}

\begin{description}
\tightlist
\item[\circled{1}]
Regresa \texttt{dict} si \texttt{d1} es un diccionario.
\item[\circled{2}]
Regresa \texttt{True} si \texttt{d1} es un diccionario, en caso
contrario regresa \texttt{False}.
\item[\circled{3}]
Regresa \texttt{True} si el valor correspondiente a \texttt{clave2} es
un entero, en caso contrario regresa \texttt{False}.
\end{description}

\begin{verbatim}
<class 'dict'>
True
False
\end{verbatim}
\end{frame}

\begin{frame}[fragile]{Acceso a claves y valores de un diccionario}
\phantomsection\label{acceso-a-claves-y-valores-de-un-diccionario}
\index{Diccionarios!Acceso a claves y valores}

Para verificar si una clave está en el diccionario \texttt{d} usamos
\texttt{clave\ in\ d}, lo cual regresará un valor lógico, \texttt{True}
o \texttt{False} dependiendo de si la clave está o no en el diccionario.
Para recuperar el valor asociado a la clave de un diccionario usamos
\texttt{d{[}\textquotesingle{}clave\textquotesingle{}{]}}, veamos

\phantomsection\label{annotated-cell-8}%
\begin{Shaded}
\begin{Highlighting}[]
\CommentTok{\textquotesingle{}Mary\textquotesingle{}} \KeywordTok{in}\NormalTok{ edad }\hspace*{\fill}\NormalTok{\circled{1}}
\NormalTok{edad[}\StringTok{"Jaime"}\NormalTok{] }\hspace*{\fill}\NormalTok{\circled{2}}
\end{Highlighting}
\end{Shaded}

\begin{description}
\tightlist
\item[\circled{1}]
Regresa \texttt{True} si la clave \texttt{Mary} está en el diccionario
\texttt{edad}, en caso contrario, regresa \texttt{False}.
\item[\circled{2}]
Regresa el valor asociado a la clave \texttt{Jaime}
\end{description}

\begin{verbatim}
False
21
\end{verbatim}
\end{frame}

\begin{frame}[fragile]{Diccionarios: \texttt{KeyError}}
\phantomsection\label{diccionarios-keyerror}
Si usas la forma
\texttt{d{[}\textquotesingle{}nombre\textquotesingle{}{]}} y
\texttt{nombre} no es una clave del diccionario, ocurre un
\texttt{KeyError}. Para evitarlo, podemos usar la sentencia \texttt{if}
para tomar decisiones, por ejemplo: en caso que se tenga el dato de
\texttt{Mary} en el diccionario se imprime su edad, de lo contrario, se
imprime un mensaje informando que ese nombre no se encuentra

\phantomsection\label{annotated-cell-9}%
\begin{Shaded}
\begin{Highlighting}[]
\NormalTok{nombre}\OperatorTok{=}\StringTok{\textquotesingle{}Mary\textquotesingle{}} \hspace*{\fill}\NormalTok{\circled{1}}
\ControlFlowTok{if}\NormalTok{ nombre }\KeywordTok{in}\NormalTok{ edad: }\hspace*{\fill}\NormalTok{\circled{2}}
  \BuiltInTok{print}\NormalTok{(}\StringTok{"La edad de }\SpecialCharTok{\%s}\StringTok{ es }\SpecialCharTok{\%g}\StringTok{"} \OperatorTok{\%}\NormalTok{ (nombre,edad[nombre]) ) }\hspace*{\fill}\NormalTok{\circled{3}}
\ControlFlowTok{else}\NormalTok{: }\hspace*{\fill}\NormalTok{\circled{4}}
  \BuiltInTok{print}\NormalTok{(}\StringTok{"No hay datos para }\SpecialCharTok{\%s}\StringTok{"} \OperatorTok{\%}\NormalTok{ nombre )}
\end{Highlighting}
\end{Shaded}

\begin{description}
\tightlist
\item[\circled{1}]
Se asigna \texttt{Mary} a nombre.
\item[\circled{2}]
Si \texttt{Mary} es una clave del diccionario \texttt{edad}.
\item[\circled{3}]
Se imprime la clave y el valor.
\item[\circled{4}]
En caso contrario, se imprime que no hay valor para esa clave.
\end{description}

\begin{verbatim}
No hay datos para Mary
\end{verbatim}
\end{frame}

\begin{frame}[fragile]{Diccionarios control \texttt{KeyError} con
\texttt{try:except}}
\phantomsection\label{diccionarios-control-keyerror-con-tryexcept}
\begin{tcolorbox}[enhanced jigsaw, leftrule=.75mm, titlerule=0mm, bottomrule=.15mm, opacitybacktitle=0.6, colbacktitle=quarto-callout-tip-color!10!white, arc=.35mm, colframe=quarto-callout-tip-color-frame, bottomtitle=1mm, title=\textcolor{quarto-callout-tip-color}{\faLightbulb}\hspace{0.5em}{Tip}, coltitle=black, colback=white, toptitle=1mm, rightrule=.15mm, breakable, toprule=.15mm, left=2mm, opacityback=0]

Como se explicó en la clase de listas, otra forma de tratar con este
error y evitar que el código \texttt{reviente} es mediante
\texttt{try:except}

\begin{Shaded}
\begin{Highlighting}[]
\NormalTok{clave}\OperatorTok{=}\StringTok{\textquotesingle{}Luis\textquotesingle{}}
\ControlFlowTok{try}\NormalTok{:}
\NormalTok{    e}\OperatorTok{=}\NormalTok{ edad[clave]}
    \BuiltInTok{print}\NormalTok{(}\SpecialStringTok{f"La edad de }\SpecialCharTok{\{}\NormalTok{clave}\SpecialCharTok{\}}\SpecialStringTok{ es }\SpecialCharTok{\{}\NormalTok{e}\SpecialCharTok{\}}\SpecialStringTok{ años"}\NormalTok{) }
\ControlFlowTok{except} \PreprocessorTok{KeyError}\NormalTok{:}
    \BuiltInTok{print}\NormalTok{(}\SpecialStringTok{f"}\SpecialCharTok{\{}\NormalTok{clave}\SpecialCharTok{\}}\SpecialStringTok{ no es una clave de edad"}\NormalTok{)}
\end{Highlighting}
\end{Shaded}

\begin{verbatim}
La edad de Luis es 15 años
\end{verbatim}

\end{tcolorbox}
\end{frame}

\section{Claves y valores a listas}\label{claves-y-valores-a-listas}

\begin{frame}[fragile]{Extraer claves y valores en listas}
\phantomsection\label{extraer-claves-y-valores-en-listas}
Tanto las claves como los valores en un diccionario se pueden extraer y
convertir a una lista, veamos \index{Diccionarios!Método!keys}
\index{Diccionarios!Método!values}

\phantomsection\label{annotated-cell-10}%
\begin{Shaded}
\begin{Highlighting}[]
\NormalTok{claves}\OperatorTok{=}\NormalTok{edad.keys() }\hspace*{\fill}\NormalTok{\circled{1}}
\BuiltInTok{print}\NormalTok{(}\BuiltInTok{type}\NormalTok{(claves)) }\hspace*{\fill}\NormalTok{\circled{2}}
\BuiltInTok{list}\NormalTok{(claves) }\hspace*{\fill}\NormalTok{\circled{3}}
\NormalTok{valores}\OperatorTok{=}\BuiltInTok{list}\NormalTok{(edad.values()) }\hspace*{\fill}\NormalTok{\circled{4}}
\BuiltInTok{print}\NormalTok{(valores) }\hspace*{\fill}\NormalTok{\circled{5}}
\end{Highlighting}
\end{Shaded}

\begin{description}
\tightlist
\item[\circled{1}]
Se recuperan las claves con el método \texttt{keys}.
\item[\circled{2}]
Note que \texttt{claves} es de tipo \texttt{dict\_keys}.
\item[\circled{3}]
Convertimos \texttt{claves} a una lista.
\item[\circled{4}]
Recuperamos los valores con el método \texttt{values} que es de tipo
\texttt{dict\_values} y lo convertimos a lista.
\item[\circled{5}]
Imprimimos los valores.
\end{description}

\begin{verbatim}
<class 'dict_keys'>
['Luis', 'Brenda', 'Diego', 'Jaime']
[15, 22, 19, 21]
\end{verbatim}
\end{frame}

\begin{frame}[fragile]{Extraer claves y valores en listas}
\phantomsection\label{extraer-claves-y-valores-en-listas-1}
\begin{itemize}
\item
  Note que \texttt{claves}, el objeto regresado por el método
  \texttt{.keys()}, no es una lista, pero podemos recuperar las claves
  como lista con la función \texttt{list()}.
\item
  Análogamente con los valores del diccionario, usando el método
  \texttt{.values()} se obtiene un objeto de tipo \texttt{dict\_values}
  que se puede convertir a lista, con la función \texttt{list()}.
\end{itemize}
\end{frame}

\section{Iteración sobre
diccionarios}\label{iteraciuxf3n-sobre-diccionarios}

\begin{frame}[fragile]{Iterar sobre un diccionario}
\phantomsection\label{iterar-sobre-un-diccionario}
De manera análoga a tuplas, listas y arreglos es posible iterar sobre
las claves de un diccionario usando bucles, por ejemplo

\begin{Shaded}
\begin{Highlighting}[]
\ControlFlowTok{for}\NormalTok{ nombre }\KeywordTok{in}\NormalTok{ edad:}
  \BuiltInTok{print}\NormalTok{(}\SpecialStringTok{f"La edad de }\SpecialCharTok{\{}\NormalTok{nombre}\SpecialCharTok{\}}\SpecialStringTok{ es }\SpecialCharTok{\{}\NormalTok{edad[nombre]}\SpecialCharTok{\}}\SpecialStringTok{ años"}\NormalTok{)}
\end{Highlighting}
\end{Shaded}

\begin{verbatim}
La edad de Luis es 15 años
La edad de Brenda es 22 años
La edad de Diego es 19 años
La edad de Jaime es 21 años
\end{verbatim}
\end{frame}

\begin{frame}[fragile]{Iterar sobre un diccionario (II)}
\phantomsection\label{iterar-sobre-un-diccionario-ii}
Suponga que deseamos la salida anterior pero en orden alfabético, por
los nombres, eso podemos lograrlo con la función \texttt{sorted()}, la
cual regresa una lista con las claves ordenadas

\phantomsection\label{annotated-cell-12}%
\begin{Shaded}
\begin{Highlighting}[]
\BuiltInTok{sorted}\NormalTok{(edad) }\hspace*{\fill}\NormalTok{\circled{1}}
\ControlFlowTok{for}\NormalTok{ nombre }\KeywordTok{in} \BuiltInTok{sorted}\NormalTok{(edad): }\hspace*{\fill}\NormalTok{\circled{2}}
  \BuiltInTok{print}\NormalTok{(}\SpecialStringTok{f"La Edad de }\SpecialCharTok{\{}\NormalTok{nombre}\SpecialCharTok{\}}\SpecialStringTok{ es }\SpecialCharTok{\{}\NormalTok{edad[nombre]}\SpecialCharTok{\}}\SpecialStringTok{ años"}\NormalTok{) }
\end{Highlighting}
\end{Shaded}

\begin{description}
\tightlist
\item[\circled{1}]
Lista con los nombres (claves del diccionario), en orden alfabético.
\item[\circled{2}]
Se recorren las claves, ordenadas alfabéticamente.
\end{description}

\begin{verbatim}
['Brenda', 'Diego', 'Jaime', 'Luis']
La Edad de Brenda es 22 años
La Edad de Diego es 19 años
La Edad de Jaime es 21 años
La Edad de Luis es 15 años
\end{verbatim}
\end{frame}

\begin{frame}[fragile]{Iterar sobre un diccionario}
\phantomsection\label{iterar-sobre-un-diccionario-1}
Tenga en cuenta que la función \texttt{sorted()} no ordena ni modifica
el diccionario, solo regresa una listas que contiene las claves
ordenadas. En el siguiente bloque de código creamos un nuevo diccionario
ordenado a partir de diccionario \texttt{edad}

\phantomsection\label{annotated-cell-13}%
\begin{Shaded}
\begin{Highlighting}[]
\NormalTok{edad\_o}\OperatorTok{=}\NormalTok{\{\} }\hspace*{\fill}\NormalTok{\circled{1}}
\ControlFlowTok{for}\NormalTok{ clave }\KeywordTok{in} \BuiltInTok{sorted}\NormalTok{(edad): }
\NormalTok{    edad\_o[clave]}\OperatorTok{=}\NormalTok{edad[clave] }\hspace*{\fill}\NormalTok{\circled{2}}
\BuiltInTok{print}\NormalTok{(edad\_o) }\hspace*{\fill}\NormalTok{\circled{3}}
\end{Highlighting}
\end{Shaded}

\begin{description}
\tightlist
\item[\circled{1}]
Creamos un diccionario vacío.
\item[\circled{2}]
Llenamos el diccionario vacío con las claves ordenadas previamente.
\item[\circled{3}]
Imprimimos.
\end{description}

\begin{verbatim}
{'Brenda': 22, 'Diego': 19, 'Jaime': 21, 'Luis': 15}
\end{verbatim}
\end{frame}

\begin{frame}[fragile]{Iterar sobre claves y valores: \texttt{.items()}}
\phantomsection\label{iterar-sobre-claves-y-valores-.items}
El método \texttt{.items()} devuelve una vista de los pares clave-valor
del diccionario como tuplas. Esto es extremadamente útil para
desempaquetar ambos elementos directamente en el bucle \texttt{for}:

\begin{Shaded}
\begin{Highlighting}[]
\ControlFlowTok{for}\NormalTok{ nombre, edad\_persona }\KeywordTok{in}\NormalTok{ edad.items():}
    \BuiltInTok{print}\NormalTok{(}\SpecialStringTok{f"La edad de }\SpecialCharTok{\{}\NormalTok{nombre}\SpecialCharTok{\}}\SpecialStringTok{ es }\SpecialCharTok{\{}\NormalTok{edad\_persona}\SpecialCharTok{\}}\SpecialStringTok{ años"}\NormalTok{)}
\end{Highlighting}
\end{Shaded}

\begin{verbatim}
La edad de Luis es 15 años
La edad de Brenda es 22 años
La edad de Diego es 19 años
La edad de Jaime es 21 años
\end{verbatim}
\end{frame}

\begin{frame}[fragile]{Ordenar elementos de un diccionario}
\phantomsection\label{ordenar-elementos-de-un-diccionario}
Alternativamente, podemos ordenar los elementos del diccionario por las
claves usando el método \texttt{.OrderedDict()} del módulo
\texttt{collections}.

\begin{Shaded}
\begin{Highlighting}[]
\ImportTok{import}\NormalTok{ collections  }
\NormalTok{EdadOr}\OperatorTok{=}\NormalTok{collections.OrderedDict(}\BuiltInTok{sorted}\NormalTok{(edad.items()))}
\BuiltInTok{print}\NormalTok{(}\BuiltInTok{dict}\NormalTok{(EdadOr))}
\end{Highlighting}
\end{Shaded}

\begin{verbatim}
{'Brenda': 22, 'Diego': 19, 'Jaime': 21, 'Luis': 15}
\end{verbatim}
\end{frame}

\begin{frame}[fragile]{Ordenar elementos de un diccionario}
\phantomsection\label{ordenar-elementos-de-un-diccionario-1}
Observe que el ordenamiento se hace de manera ascendente, para hacerlo
en orden descendente use \texttt{reverse=True} en la función
\texttt{sorted()}
\end{frame}

\begin{frame}[fragile]{Hacer copia de un diccionario}
\phantomsection\label{hacer-copia-de-un-diccionario}
Se puede hacer una copia de un diccionario mediante el método
\texttt{.copy()} y como en el caso de las listas y arreglos, se puede
borrar una entrada de un diccionario con \texttt{del}

\phantomsection\label{annotated-cell-16}%
\begin{Shaded}
\begin{Highlighting}[]
\NormalTok{edadC}\OperatorTok{=}\NormalTok{edad.copy() }\hspace*{\fill}\NormalTok{\circled{1}}
\KeywordTok{del}\NormalTok{ edadC[}\StringTok{"Luis"}\NormalTok{] }\hspace*{\fill}\NormalTok{\circled{2}}
\NormalTok{edadC}
\end{Highlighting}
\end{Shaded}

\begin{description}
\tightlist
\item[\circled{1}]
Copia del diccionario edad.
\item[\circled{2}]
Se borra de \texttt{edadC} la clave \texttt{Luis} junto con su valor.
\end{description}

\begin{verbatim}
{'Brenda': 22, 'Diego': 19, 'Jaime': 21}
\end{verbatim}
\end{frame}

\section{Métodos aplicables a
diccionarios}\label{muxe9todos-aplicables-a-diccionarios}

\begin{frame}[fragile]{Método \texttt{.pop()}}
\phantomsection\label{muxe9todo-.pop}
La mayoría de los métodos que estudiamos para listas son aplicables a
diccionarios. \index{Diccionarios!Método!\texttt{.pop()}}

El \emph{método \texttt{pop()}} remueve específicamente una clave de
diccionario y devuelve valor correspondiente. En caso que la clave dada
como argumento no se encuentre en el diccionario, \texttt{pop()} lanza
una excepción \texttt{KeyError}.

\phantomsection\label{annotated-cell-17}%
\begin{Shaded}
\begin{Highlighting}[]
\NormalTok{edadC}\OperatorTok{=}\NormalTok{edad.copy() }\hspace*{\fill}\NormalTok{\circled{1}}
\NormalTok{edadC.pop(}\StringTok{"Luis"}\NormalTok{) }\hspace*{\fill}\NormalTok{\circled{2}}
\NormalTok{edadC}
\end{Highlighting}
\end{Shaded}

\begin{description}
\tightlist
\item[\circled{1}]
Copia del diccionario edad.
\item[\circled{2}]
Regresa la edad de Luis y lo elimina del diccionario.
\end{description}

\begin{verbatim}
15
{'Brenda': 22, 'Diego': 19, 'Jaime': 21}
\end{verbatim}
\end{frame}

\begin{frame}[fragile]{Método \texttt{.popitem()}}
\phantomsection\label{muxe9todo-.popitem}
\index{Diccionarios!Método!\texttt{.popitem()}}

Observe que en \texttt{edadC} ya no está la clave \texttt{Luis}.

El \emph{método \texttt{.popitem()}} funciona como \texttt{pop()}, pero
con la última clave, es decir, elimina el último ítem del diccionario y
regresa una tupla con la clave y el valor correspondientes.

\begin{Shaded}
\begin{Highlighting}[]
\NormalTok{edadC}\OperatorTok{=}\NormalTok{edad.copy() }\CommentTok{\# copia de edad }
\NormalTok{edadC.popitem() }
\end{Highlighting}
\end{Shaded}

\begin{verbatim}
('Jaime', 21)
\end{verbatim}
\end{frame}

\begin{frame}[fragile]{Método \texttt{update()}}
\phantomsection\label{muxe9todo-update}
\index{Diccionarios!Método!\texttt{.update()}}

El \emph{método \texttt{update}} permite actualizar la información
contenida en un diccionario. Recibe como parámetro otro diccionario: si
los dos diccionarios tienen claves en común, actualiza los valores de
las claves comunes; si no hay claves comunes, los pares
\texttt{clave:valor} no comunes serán agregados al diccionario.

\phantomsection\label{annotated-cell-19}%
\begin{Shaded}
\begin{Highlighting}[]
\NormalTok{edad2}\OperatorTok{=}\NormalTok{\{}\StringTok{"Brenda"}\NormalTok{:}\DecValTok{25}\NormalTok{,}\StringTok{"Mary"}\NormalTok{:}\DecValTok{18}\NormalTok{\} }\hspace*{\fill}\NormalTok{\circled{1}}
\NormalTok{edad.update(edad2) }\hspace*{\fill}\NormalTok{\circled{2}}
\NormalTok{edad }
\end{Highlighting}
\end{Shaded}

\begin{description}
\tightlist
\item[\circled{1}]
Nuevo diccionario, note que Brenda ya está en el diccionario
\texttt{edad}
\item[\circled{2}]
Se actualiza el diccionario \texttt{edad}, observe que a Brenda, que es
una clave común a los dos diccionarios, se le actualiza la edad.
\end{description}

\begin{verbatim}
{'Luis': 15, 'Brenda': 25, 'Diego': 19, 'Jaime': 21, 'Mary': 18}
\end{verbatim}
\end{frame}

\begin{frame}[fragile]{Método \texttt{.get()}}
\phantomsection\label{muxe9todo-.get}
\index{Diccionarios!Método!\texttt{.get()}} El \emph{método
\texttt{get()}} recibe una clave y devuelve el valor que le corresponde
en el diccionario.\\
\index{Diccionarios!Método!\texttt{.setdefault()}} El \emph{método
\texttt{.setdefault()}} tiene dos usos: el primero es como el método
\texttt{get()} y, el segundo, para agregar un nuevo elemento al
diccionario

\phantomsection\label{annotated-cell-20}%
\begin{Shaded}
\begin{Highlighting}[]
\NormalTok{edad.get(}\StringTok{"Brenda"}\NormalTok{) }\hspace*{\fill}\NormalTok{\circled{1}}
\NormalTok{edad.setdefault(}\StringTok{"Brenda"}\NormalTok{) }\hspace*{\fill}\NormalTok{\circled{2}}
\NormalTok{edad.setdefault(}\StringTok{"Lila"}\NormalTok{,}\DecValTok{16}\NormalTok{) }\hspace*{\fill}\NormalTok{\circled{3}}
\ControlFlowTok{for}\NormalTok{ clave }\KeywordTok{in}\NormalTok{ edad: }
    \BuiltInTok{print}\NormalTok{(}\SpecialStringTok{f"\textquotesingle{}}\SpecialCharTok{\{}\NormalTok{clave}\SpecialCharTok{\}}\SpecialStringTok{\textquotesingle{}:}\SpecialCharTok{\{}\NormalTok{edad[clave]}\SpecialCharTok{\}}\SpecialStringTok{"}\NormalTok{) }\hspace*{\fill}\NormalTok{\circled{4}}
\end{Highlighting}
\end{Shaded}

\begin{description}
\tightlist
\item[\circled{1}]
\texttt{get()}: recupera el valor de la clave Brenda.
\item[\circled{2}]
\texttt{setdefault}, primer uso: exactamente lo mismo que
\texttt{get()}.
\item[\circled{3}]
\texttt{setdefault}, segundo uso: agrega al diccionario \texttt{edad}
una nueva pareja \texttt{clave:valor}, \texttt{Lila:16}.
\item[\circled{4}]
Recorremos el diccionario para imprimir de forma vertical.
\end{description}

\begin{verbatim}
25
25
16
'Luis':15
'Brenda':25
'Diego':19
'Jaime':21
'Mary':18
'Lila':16
\end{verbatim}
\end{frame}

\begin{frame}[fragile]{Tuplas a diccionarios}
\phantomsection\label{tuplas-a-diccionarios}
Considere la siguiente tupla, conformada por 6 tuplas anidadas:

\begin{Shaded}
\begin{Highlighting}[]
\NormalTok{tup }\OperatorTok{=}\NormalTok{ ((}\StringTok{\textquotesingle{}A\textquotesingle{}}\NormalTok{, }\DecValTok{5}\NormalTok{), (}\StringTok{\textquotesingle{}A\textquotesingle{}}\NormalTok{, }\DecValTok{7}\NormalTok{),}
\NormalTok{       (}\StringTok{\textquotesingle{}B\textquotesingle{}}\NormalTok{, }\DecValTok{2}\NormalTok{),(}\StringTok{\textquotesingle{}B\textquotesingle{}}\NormalTok{, }\DecValTok{9}\NormalTok{),(}\StringTok{\textquotesingle{}B\textquotesingle{}}\NormalTok{, }\DecValTok{1}\NormalTok{),}
\NormalTok{       (}\StringTok{\textquotesingle{}C\textquotesingle{}}\NormalTok{,}\DecValTok{4}\NormalTok{) )}
\end{Highlighting}
\end{Shaded}

Note que \texttt{\textquotesingle{}A\textquotesingle{}} está dos veces;
\texttt{\textquotesingle{}B\textquotesingle{}}, 3 veces y
\texttt{\textquotesingle{}C\textquotesingle{}} una vez en la tupla. El
ejercicio consiste en trasladar esta información a un diccionario donde
\texttt{\textquotesingle{}A\textquotesingle{}},
\texttt{\textquotesingle{}B\textquotesingle{}} y
\texttt{\textquotesingle{}C\textquotesingle{}} conformen las claves y
los valores que coresponden a cada clave se organicen en una lista.
\end{frame}

\begin{frame}[fragile]{Tuplas a diccionarios (cont)}
\phantomsection\label{tuplas-a-diccionarios-cont}
\phantomsection\label{annotated-cell-22}%
\begin{Shaded}
\begin{Highlighting}[]
\NormalTok{dicn }\OperatorTok{=}\NormalTok{ \{\} }\hspace*{\fill}\NormalTok{\circled{1}}
\ControlFlowTok{for}\NormalTok{ clave, valor }\KeywordTok{in}\NormalTok{ tup: }\hspace*{\fill}\NormalTok{\circled{2}}
    \ControlFlowTok{if}\NormalTok{ clave }\KeywordTok{in}\NormalTok{ dicn: }\hspace*{\fill}\NormalTok{\circled{3}}
\NormalTok{        dicn[clave].append(valor) }\hspace*{\fill}\NormalTok{\circled{4}}
    \ControlFlowTok{else}\NormalTok{: }\hspace*{\fill}\NormalTok{\circled{5}}
\NormalTok{        dicn[clave] }\OperatorTok{=}\NormalTok{ [valor] }\hspace*{\fill}\NormalTok{\circled{6}}
\NormalTok{dicn }
\end{Highlighting}
\end{Shaded}

\begin{description}
\tightlist
\item[\circled{1}]
Diccionario vacío para guardar los datos.
\item[\circled{2}]
Se recorre la tupla.
\item[\circled{3}]
En caso que la clave ya esté en dicn.
\item[\circled{4}]
Agrega el valor a la lista
\item[\circled{5}]
En caso que la clave aún no esté
\item[\circled{6}]
Anexa a \texttt{dicn} la clave con el respectivo valor
\end{description}

\begin{verbatim}
{'A': [5, 7], 'B': [2, 9, 1], 'C': [4]}
\end{verbatim}
\end{frame}

\begin{frame}[fragile]{}
\phantomsection\label{section}
Observe que la tupla inicial tiene longitud 6, mientras que la longitud
del diccionario, que contiene la misma información es 3. Esta forma de
guardar los datos es mucho más eficiente en cuanto al uso de memoria y
espacio en disco duro. Otra forma de realizar la misma tarea es mediante
la combinación de los métodos \texttt{.setdefault()} y
\texttt{.append()}

\phantomsection\label{annotated-cell-23}%
\begin{Shaded}
\begin{Highlighting}[]
\NormalTok{dicn }\OperatorTok{=}\NormalTok{ \{\} }\hspace*{\fill}\NormalTok{\circled{1}}
\ControlFlowTok{for}\NormalTok{ clave, valor }\KeywordTok{in}\NormalTok{ tup: }\hspace*{\fill}\NormalTok{\circled{2}}
\NormalTok{    dicn.setdefault(clave, []).append(valor) }\hspace*{\fill}\NormalTok{\circled{3}}
\NormalTok{dicn }
\end{Highlighting}
\end{Shaded}

\begin{description}
\tightlist
\item[\circled{1}]
Diccionario vacío para guardar los datos.
\item[\circled{2}]
Se recorre la tupla.
\item[\circled{3}]
Se agrega la información al diccionario:
\texttt{setdefault(clave,\ {[}{]})} verifica si la clave no está la
agrega con una lista vacía como valor, luego \texttt{append(valor)}
agrega el valor a la lsita.
\end{description}

\begin{verbatim}
{'A': [5, 7], 'B': [2, 9, 1], 'C': [4]}
\end{verbatim}
\end{frame}

\begin{frame}[fragile]{Comprensión de Diccionarios}
\phantomsection\label{comprensiuxf3n-de-diccionarios}
Al igual que con las listas, podemos generar nuevos diccionarios de
forma compacta utilizando una sola línea de código.

Supongamos que necesitamos calcular el cuadrado de las edades:

\begin{Shaded}
\begin{Highlighting}[]
\NormalTok{edades\_cuadrado }\OperatorTok{=}\NormalTok{ \{nombre: valor}\OperatorTok{**}\DecValTok{2} \ControlFlowTok{for}\NormalTok{ nombre, valor }\KeywordTok{in}\NormalTok{ edad.items()\}}
\BuiltInTok{print}\NormalTok{(edades\_cuadrado)}
\end{Highlighting}
\end{Shaded}

\begin{verbatim}
{'Luis': 225, 'Brenda': 625, 'Diego': 361, 'Jaime': 441, 'Mary': 324, 'Lila': 256}
\end{verbatim}
\end{frame}

\section{Aplicación}\label{aplicaciuxf3n}

\begin{frame}[fragile]{Polinomios como diccionarios}
\phantomsection\label{polinomios-como-diccionarios}
Las claves en los diccionarios no tienen que ser necesariamente cadenas
de caracteres \emph{(strings)}, en realidad cualquier objeto
\textbf{inmutable} (como las tuplas) de Python puede usarse como clave.
En este ejemplo se muestra como los diccionarios con enteros como claves
se pueden usar para representar polinomios. Los datos asociados con el
polinomio
\begin{equation}\phantomsection\label{eq-polinomio}{p(x)=-1+x^2+3x^7 }\end{equation}
se pueden expresar como un conjunto de pares de la forma
\texttt{clave:valor}, donde la clave corresponde al exponente y el valor
al coeficiente
\end{frame}

\begin{frame}[fragile]{}
\phantomsection\label{section-1}
Las entradas del diccionario, de acuerdo con esto, son
\{\texttt{exponente:coeficiente}\}.

En el caso del polinomio dado en (\ref{eq-polinomio}), el coeficiente
\(-1\) corresponde al exponente cero para la variable \(x\), ya que
\(-1=-1x^0\), el coeficiente 1 pertenece al exponente 2, el coeficiente
3 corresponde al exponente 7. A continuación se usa el diccionario
\texttt{p} para guardar la información de (\ref{eq-polinomio})

\begin{Shaded}
\begin{Highlighting}[]
\NormalTok{p}\OperatorTok{=}\NormalTok{\{}\DecValTok{0}\NormalTok{:}\OperatorTok{{-}}\DecValTok{1}\NormalTok{, }\DecValTok{2}\NormalTok{:}\DecValTok{1}\NormalTok{, }\DecValTok{7}\NormalTok{:}\DecValTok{3}\NormalTok{\} }
\ControlFlowTok{for}\NormalTok{ m }\KeywordTok{in}\NormalTok{ p:}
        \BuiltInTok{print}\NormalTok{(}\SpecialStringTok{f"Coeficiente: }\SpecialCharTok{\{}\NormalTok{p[m]}\SpecialCharTok{\}}\SpecialStringTok{, Exponente: }\SpecialCharTok{\{}\NormalTok{m}\SpecialCharTok{\}}\SpecialStringTok{"}\NormalTok{)  }
\end{Highlighting}
\end{Shaded}

\begin{verbatim}
Coeficiente: -1, Exponente: 0
Coeficiente: 1, Exponente: 2
Coeficiente: 3, Exponente: 7
\end{verbatim}
\end{frame}

\begin{frame}[fragile]{}
\phantomsection\label{section-2}
Podemos definir una función que recibe un polinomio (expresado como un
diccionario) y un valor de la variable \(x_0\) y evalúe \(p(x_0)\)

\phantomsection\label{annotated-cell-26}%
\begin{Shaded}
\begin{Highlighting}[]
\KeywordTok{def}\NormalTok{ evalpoli(poli,x): }\hspace*{\fill}\NormalTok{\circled{1}}
    \ControlFlowTok{return}  \BuiltInTok{sum}\NormalTok{([poli[m]}\OperatorTok{*}\NormalTok{x}\OperatorTok{**}\NormalTok{m }\ControlFlowTok{for}\NormalTok{ m }\KeywordTok{in}\NormalTok{ poli]) }\hspace*{\fill}\NormalTok{\circled{2}}
\NormalTok{evalpoli(p,}\DecValTok{1}\NormalTok{) }\hspace*{\fill}\NormalTok{\circled{3}}
\end{Highlighting}
\end{Shaded}

\begin{description}
\tightlist
\item[\circled{1}]
Definimos la función.
\item[\circled{2}]
Observe el uso de comprensión de lista. Se recorren las claves del
diccionario y se multiplica \texttt{valor} por \texttt{x} elevado a
\texttt{clave}.
\item[\circled{3}]
Se hace el llamado a la función \texttt{evalpoli}, entregándole el
polinomio \texttt{p} y el valor donde se avaluará.
\end{description}

\begin{verbatim}
3
\end{verbatim}
\end{frame}

\begin{frame}[fragile]{}
\phantomsection\label{section-3}
Así que el polinomio \ref{eq-polinomio} en \(x_0=1\), toma un valor de
3, es decir, \(p(1)=3\).

El la parte esencial de la función \texttt{evalpoli} es que la variable
\texttt{m} varía sobre las claves del diccionario, luego se recupera el
valor correspondiente a esa clave con \texttt{poli{[}m{]}}, y se
multiplica por \texttt{x} elevado a la \texttt{m}.
\end{frame}




\end{document}
